\documentclass[12pt,a4paper]{article}
\usepackage[brazil]{babel}
\usepackage{geometry}
\usepackage{hyperref}
\usepackage{tikz}
\usetikzlibrary{positioning}

\geometry{margin=2.5cm}

\title{Projeto Final\\Estrutura de Dados e Complexidade de Algoritmos\\TSP com Heurística e VND}
\author{Joaquim Breno Brito Cavalcante \\ Universidade Federal da Paraíba \\ Centro de Informática \\ Prof. Gilberto Farias}
\date{Junho de 2026}

\begin{document}
\maketitle

\section{Definição do POC}

Escolhi o Problema do Caixeiro Viajante, o famoso TSP.

A ideia é simples, tem um monte de cidades e o caixeiro precisa visitar todas, uma vez cada, e voltar pra onde começou. O objetivo é gastar o mínimo de distância (ou custo) possível no caminho todo.

Dá pra pensar num grafo completo onde cada par de cidades tem um custo de ir de uma pra outra, geralmente a distância. A solução é só a ordem em que as cidades são visitadas, como ir de A pra C, depois B, depois D, e fechar o ciclo.

Regras:
\begin{itemize}
    \item passa em cada cidade exatamente uma vez;
    \item no final volta pra cidade inicial.
\end{itemize}


\section{Prova de pertencimento a NP}

Na aula de NP-completude a gente viu que primeiro convém olhar a versão de decisão do problema.

No TSP-DEC a pergunta é: dado os custos entre cidades e um número k, existe alguma rota que visita tudo e custa no máximo k?

\textbf{Certificado:} a lista com a ordem das cidades. Tipo: 3, 1, 5, 2, 4.

\textbf{Como verificar:} alguém pega essa lista e:
\begin{enumerate}
    \item vê se todas as cidades aparecem uma vez só, sem repetir e sem faltar ninguém.
    \item soma o custo de cada trecho, incluindo a volta pro início.
    \item responde sim se der menor ou igual a k, senão responde não.
\end{enumerate}

Isso dá pra fazer em tempo linear no número de cidades. Então o TSP-DEC está em NP.

\section{NP-difícil e redução}

Pra mostrar que é difícil de verdade, reduz o Ciclo Hamiltoniano (CH) pro TSP-DEC. CH é aquele problema clássico: existe um ciclo que passa por todos os vértices do grafo uma vez? Isso é NP-completo, apareceu nos slides da disciplina.

\textbf{Como montar a redução}

Pega a instância do CH e transforma assim:
\begin{itemize}
    \item mesmas cidades (vértices);
    \item se existe aresta entre duas cidades, o custo vira 1;
    \item se não existe aresta, o custo vira 2;
    \item k fica igual ao total de cidades.
\end{itemize}

Montar a matriz de custos leva tempo quadrático no número de cidades, o que ainda é polinomial.

\textbf{Por que funciona}

Se o grafo original tem ciclo hamiltoniano, dá pra montar um tour no TSP só com arestas de custo 1. Aí o total fica igual ao número de cidades, e isso passa no limite k.

Se não tem ciclo hamiltoniano, qualquer tour completo vai ser obrigado a usar pelo menos uma aresta que não existia no grafo original, ou seja, custo 2. O melhor caso vira número de cidades mais 1, que é maior que k.

Com isso CH reduz polinomialmente pro TSP-DEC. CH é NP-completo, então TSP-DEC é NP-difícil. E como TSP-DEC também está em NP, ele é NP-completo. A versão de otimização (achar o menor custo) herda a mesma dificuldade.

\begin{center}
\begin{tikzpicture}[node distance=2.8cm, >=stealth, ->]
  \node (ch) {Ciclo Hamiltoniano};
  \node (tspd) [right=of ch] {TSP-DEC};
  \node (tspo) [right=of tspd] {TSP otimização};
  \draw (ch) -- node[above, font=\small] {custo 1 se tem aresta} (tspd);
  \draw (tspd) -- (tspo);
\end{tikzpicture}
\end{center}

\section{Instâncias da literatura}

Usei instâncias do TSPLIB, que é o repositório clássico de benchmark pra TSP: \url{http://comopt.ifi.uni-heidelberg.de/software/TSPLIB95/}

\begin{table}[h]
\centering
\begin{tabular}{lrrl}
\hline
Instância & Cidades & Ótimo & Observação \\
\hline
eil51     & 51  & 426   & pequena, clássica \\
berlin52  & 52  & 7542  & coordenadas euclidianas \\
st70      & 70  & 675   & tamanho médio \\
kroA100   & 100 & 21282 & a maior que testei \\
\hline
\end{tabular}
\caption{Instâncias TSPLIB usadas nos testes.}
\end{table}

\section{Implementação e resultados (A2)}

Código no GitHub:
\url{https://github.com/JoaquimBreno/atividades_modulo3_estrutura}

\subsection{Representação da solução}

Guardei a solução como uma lista de inteiros, cada um é uma cidade na ordem da visita. Exemplo: [0, 3, 1, 4, 2].

O custo é a soma das distâncias de cada par consecutivo, mais o trecho de volta da última cidade pra primeira.

Na leitura do arquivo .tsp eu monto a matriz de distâncias inteira e deixo na memória, assim não fico recalculando coordenada o tempo todo.

\subsection{Heurística de construção: vizinho mais próximo}

Começa numa cidade (sorteada aleatoriamente nos testes) e, enquanto sobrar cidade, vai sempre pra mais perto que ainda não visitou. Quando acaba, o ciclo fecha sozinho na hora de calcular o custo.

Roda em tempo quadrático no número de cidades. Simples e rápido de implementar, como vi na aula.

\subsection{Movimentos de vizinhança}

Usei três vizinhanças no VND:
\begin{itemize}
    \item 2-opt: inverte um pedaço da rota;
    \item swap: troca duas cidades de posição;
    \item relocate: tira uma cidade de um lugar e encaixa em outro.
\end{itemize}

\subsection{VND}

O VND pega a solução do vizinho mais próximo e tenta melhorar. Passa pelas vizinhanças nessa ordem: 2-opt, swap, relocate.

Se alguma melhora, volta pro 2-opt. Se nenhuma das três consegue melhorar, para. Ótimo local, sem frescura.

\subsection{Resultados computacionais}

Rodei cada instância 10 vezes com cidade inicial aleatória. A tabela mostra o melhor custo que apareceu, o ótimo conhecido da literatura, o gap em porcento e o tempo médio.

\begin{table}[h]
\centering
\begin{tabular}{lrrrr}
\hline
Instância & Melhor & Ótimo & Gap (\%) & Tempo médio (s) \\
\hline
berlin52 & 7542 & 7542 & 0.0 & 0.035 \\
eil51    & 429  & 426  & 0.7 & 0.047 \\
kroA100  & 21379 & 21282 & 0.5 & 0.362 \\
st70     & 683  & 675  & 1.2 & 0.151 \\
\hline
\end{tabular}
\caption{Resultados do VND nas instâncias TSPLIB.}
\end{table}

Os gaps ficaram abaixo de 1,5\% em todos os casos. Pra uma heurística simples, sem tabu search nem simulated annealing, acho que tá ok. O berlin52 chegou no ótimo em uma das execuções. O tempo continua baixo, frações de segundo até com 100 cidades.

\end{document}
